\subsection{Đau lưng cấp}

\begin{longtblr}[
    caption = {Đau lưng cấp},
]{
    width = \textwidth,
    colspec = {| X[1,c] |X[1, c]|X[2, c]|X[1, c]|X[7, l]|},
    % hlines,
    rowhead = 1,
}
    \hline
    ID & Động tác & Vùng tác động & Số lần & Huyệt
    \\ \hline
    1 & Day & Hai bên cột sống từ D7 đến mông & 3 &
    \\ \hline
    2 & Lăn & Hai bên cột sống từ D7 đến mông & 3 &
    \\ \hline
    3 & Bóp & Hai bên cột sống từ D7 đến mông & 3 &
    \\ \hline
    4 & Bấm & Huyệt & 1 & {
        Giáp Tích nơi đau, Thận Du, Đại Trường Du, Cách Du, a thị huyệt.
    }
    \\ & & & & {
        Giáp Tích

        \giaptich
    }
    \\ & & & & {
        Thận Du 
        
        \thandu
    }
    \\ & & & & {
        Đại Trường Du

        \daitruongdu
    }
    \\ & & & & {
        Cách Du       

        \cachdu
    }
    \\ \hline
    5 & Nắn & Cột sống & 1 & Nếu đốt sống bị lệch
    \\ \hline
    6 & Vận động & Cột sống & 1 & {
    Bệnh nhân nằm nghiêng chân dưới duỗi thẳng, chân trên co, tay ở phía dưới để trước mặt, tay ở phía trên để quặt sau lưng. Một cẳng tay thầy thuốc để ở rãnh delta ngực, một cẳng tay để ở mông, hai tay vận động ngược chiều nhau một cách nhẹ nhàng, sau đó đột nhiên làm mạnh một cái sẽ phát ra tiếng kêu khục, rồi đổi chiều làm bên đối diện.
    \newline
    Gập đùi vào ngực: bệnh nhân nằm ngửa co chân, một tay thầy thuốc để ở hai đàu gối, một tay để dưới mông, hai tay phối hợp nhịp nhàng nâng mông lên và ấn gối xuống, từ từ tăng dần khi nào đầu gối sát ngực thì đưa mông sang phải và sang trái hai lần.
    \newline
    Chú ý: động tác gập đùi vào ngực nếu có nghi ngờ thoát vị đĩa đệm thì không làm.
    }
    \\ \hline
    7 & Phát & Thắt lưng & 1 &
    \\ \hline
\end{longtblr}
